\section{Introduction}
\label{sec:intro}

Online social platforms give a direct view of how communities form and exchange information. Bluesky has grown quickly as an alternative to X/Twitter and has attracted, in particular, a large and active scientific community of researchers, science journalists and outreach organizations. That makes it a useful case study: a self-selected population whose online ties reflect real scholarly and professional relationships.

This project collects, characterizes and analyses the follow network of the scientific community on Bluesky. Starting from five thematic seed accounts, a snowball crawl over the platform's public endpoints yields a connected network of 15{,}000 accounts and about 1.67M follow relationships, enriched with node attributes (profile metadata) and a sample of posts. Three guiding questions are addressed: what is the structure of this network and does it differ from classical random models? Is it modular, and do its communities correspond to recognizable scientific domains? And how do information and behaviors spread over it?

The report is organized along the parts of the assignment. Section~\ref{sec:data} describes the data collection and the resulting network; Section~\ref{sec:analysis} characterizes its topology and contrasts it with Erd\H{o}s--R\'enyi and Barab\'asi--Albert baselines; Section~\ref{sec:cd} detects and interprets its community structure; Section~\ref{sec:part3} addresses link prediction and diffusion (including a game-theoretic cascade); and Section~\ref{sec:open} tackles the open question of topical echo chambers. Together they describe a dense, small-world, approximately scale-free network with a strong community structure aligned with academic disciplines, and topical chambers that are permeable, not sealed; the final discussion draws these threads together.
